\documentclass[10pt,a4paper]{amsart}
\usepackage[margin=19mm]{geometry}
\usepackage{amsmath,amssymb,mathtools}
\usepackage[scr=rsfs,cal=euler]{mathalfa}
\usepackage{xfrac}
\usepackage[unicode,hidelinks,bookmarks=false]{hyperref}
\newcommand{\ie}{i.e.\ }
\newcommand{\cf}{cf.\ }
\newcommand{\resp}{respectively\ }
\newcommand{\Subsection}{\subsection}
\providecommand{\nobreakdash}{\nobreak\hbox{-}}
\setlength{\emergencystretch}{2em}
\multlinegap0pt
\usepackage{bm}
\usepackage{bbm}
\usepackage{dsfont}
\usepackage{xspace}
\usepackage{cancel}
\usepackage{mathtools}
\usepackage{amsthm}
\makeatletter
\@ifundefined{lemma}{\newtheorem{lemma}{Lemma}}{}
\makeatother
\title[Global Dynamics of Trait-Structured Generalised Lotka-Volter]{Global Dynamics of Trait-Structured Generalised Lotka-Volterra Systems with Trait-Independent Interactions}
\author{Manh Hong Duong, Nataliya Balabanova, Blaine van Rensburg}
\date{Source: arXiv:2607.25573v1 (CC BY 4.0)}
\begin{document}
\maketitle

\noindent\emph{Source and licence.} Global Dynamics of Trait-Structured Generalised Lotka-Volterra Systems with Trait-Independent Interactions. Version arXiv:2607.25573v1 (CC BY 4.0), distributed under the Creative Commons Attribution 4.0 International licence, as recorded in the arXiv licence metadata for this exact version and shown on the abstract page https://arxiv.org/abs/2607.25573v1. Full rights notice: licenses/arxiv-2607.25573-CC-BY-4.0.txt.\par\smallskip

\noindent\textit{Editorial scope.} Counted unit: the Lemma whose source label is \texttt{lem: dynamics LV} in the article (source0.tex lines 996--999, source label \texttt{lem: dynamics LV}), delivered together with the article's own complete original proof of it (source0.tex lines 1000--1033, one \texttt{proof} environment of 34 lines). No mathematics is written, repaired or re-derived: statement and proof are the article's own text line for line. Required context retained uncounted: eqn:PopDynFull (lines 838--843); eqn:PopDynFull\_3 (lines 850--852). Every definition, display and label of those lines is retained unchanged. Omitted from this package and not redistributed: 1--837, 844--849, 853--995, 1034--1054. Nothing inside a retained excerpt is omitted or altered: every retained body is delivered exactly as the article's lines are written, so no omission receipt of any kind accompanies this package. Citations: the retained excerpts carry no citation command at all, so no bibliography entry of the article is transcribed and no reference is redistributed. Reference adaptation: none was needed and none was made; every reference inside the delivered unit resolves inside it, so no unresolved reference or citation is produced. Printed slips kept literal: no printed word, symbol, hypothesis or number of the counted statement or of its proof was altered. Layout and class: the article's own class is \texttt{article} with packages including graphicx, xcolor, amsmath, amsfonts, subfigure, caption, url, hyperref, amsthm, biblatex; the wrapper is the lane's pinned \texttt{amsart} editorial wrapper at 10pt on A4, which restates the theorem-like environments the retained text needs and adds the article's own macro definitions that the retained text needs (none), with extra wrapper packages (bm, bbm, dsfont, xspace, cancel, mathtools). The delivered numbering is the unit's own. Assets: no retained span contains a figure environment, a graphics-inclusion command, a PDF-page inclusion, a table or a TikZ picture; no image file is copied into this package, the package declares no asset dependency and no image file is redistributed with it. Licence: version arXiv:2607.25573v1 of this article is distributed under the Creative Commons Attribution 4.0 International licence, as recorded in the arXiv licence metadata for this exact version and shown on the abstract page https://arxiv.org/abs/2607.25573v1. A full rights notice is delivered at licenses/arxiv-2607.25573-CC-BY-4.0.txt. Characters: every retained excerpt is pure ASCII, so the delivered engine, XeLaTeX with its default Latin Modern text font, reports no missing character.
\begin{equation}\label{eqn:PopDynFull}
    \begin{cases}
        \frac{d\rho_i}{dt}=\rho_i\left(-\lambda_i+\alpha_i(t)-\sum_{i=1}^{N}a_{ij}\rho_j\right),\\
        \rho_i(0)=\int_{\mathbb{R}}u_{i,0}(y)dy.
    \end{cases}
\end{equation}

\begin{equation}\label{eqn:PopDynFull_3}
        \frac{d\widetilde{\rho}_i}{dt}=\widetilde{\rho}_i\left(-\lambda_i-\sum_{i=1}^{N}a_{ij}\widetilde{\rho}_j\right),
   \end{equation}

\begin{lemma}
\label{lem: dynamics LV}
    Assume that the trajectory of an unperturbed equation \eqref{eqn:PopDynFull_3} has a periodic cycle $C$ of width $diam (C)$  and that $M_1$ is the maximum value of $||\rho(t)||$ on this trajectory. Then, if the time delay $\tau$  was chosen to be sufficiently large, the time $T'$ that the perturbed trajectory will spend in the $2\epsilon$ vicinity of the unperturbed one will be $t\le \min\left\{\frac{\epsilon}{\mathrm{diam} C}, \frac{\epsilon}{M_1\max_{\xi\ge\tau}||\alpha(\xi)||}\right\}$.
\end{lemma}

\begin{proof}
For brevity, we stick to the notation of the theorem above, and denote the flow of the perturbed system by $\lambda(t,\tau,x)$ and the unperturbed by $\phi(t,x)$, with the left hand side of \eqref{eqn:PopDynFull} denoted by $f(t,x)$ and the right hand side of the unperturbed equation by $g(x)$.

When the $\omega$-limit set of the unperturbed trajectory is a point, it is clear that the two solutions will stay in the $\epsilon-$vicinity of each other indefinitely long, provided they started sufficiently close. In the case when the $\omega$-limit set is a cycle, things are a bit different. 

Using the definition of the flow, we can write the following estimations for the distance between the solutions:
\begin{equation*}
    \begin{split}
        ||\lambda(t + \tau, \tau,x) - \phi(t,x)||&\le \int\limits_0^t ||f(\lambda(s + \tau,s,x)) -g(\phi(s,x))||\,d s\\&\le\int\limits_0^t ||f(\lambda(s + \tau,s,x)) -g(\lambda(s + \tau,s,x))||\,d s \\&+ \int\limits_0^t ||g(\lambda(s + \tau,s,x)) -g(\phi(s,x))||\, d s.
    \end{split}
\end{equation*}
The first summand on the right hand side is proportional to $\alpha:$
\begin{equation*}
    \begin{split}
       \int\limits_0^t ||f(\lambda(s + \tau,\tau,x)) -g(\lambda(s + \tau,\tau,x))||d s & = \int\limits_0^t||\alpha(s + \tau)\rho(\lambda(s + \tau,\tau,x){)} ||\, d s.
    \end{split}
\end{equation*}
Here, the integrand is a vector with entries $(\alpha_i {\rho}_i)_i$ and the norm is the standard Euclidean vector norm.  Assume that $||{\rho}_i(t)||<M_1$ for some constant $M_1$ -- this follows from the $\omega$-limit set being the periodic trajectory. Then 
\[
\int\limits_0^t ||f(\lambda(s + \tau,s,x)) -g(\lambda(s + \tau,s,x))||d s \le M_1 t \max_{\xi\ge \tau}||\alpha(\xi)||. 
\]
For the second part, we have 
\begin{equation*}
    \begin{split}
       \int\limits_0^t ||g(\lambda(s + \tau,s,x)) -g(\phi(s,x))||d s&\le \mathrm{diam}\ C*t ,
    \end{split}
\end{equation*}
where $\mathrm{diam}\ C$ is the diameter of the periodic trajectory. 

Therefore, if we want the perturbed trajectory to stay in the $2\epsilon$-neighbourhood of the unperturbed one, it must hold that 
\[
t\le \min\left\{\frac{\epsilon}{\mathrm{diam} C}, \frac{\epsilon}{M_1\max_{\xi\ge\tau}||\alpha(\xi)||}\right\}.
\]
\end{proof}

\end{document}
